\section{Data Collection} \label{sec:collect}

% \todo{Cite ParlAI somewhere?}

The main goal of our work is to collect a diverse and explainable question answering dataset that requires multi-hop reasoning.
% As discussed in the previous section,
One way to do so is to define reasoning chains based on a knowledge base \cite{welbl2017constructing,talmor2018web}.
However, the resulting datasets are limited by the incompleteness of entity relations and the lack of diversity in the question types.
Instead, in this work, we focus on text-based question answering in order to diversify the questions and answers.
The overall setting is that given some context paragraphs (e.g., a few paragraphs, or the entire Web) and a question, a QA system answers the question by extracting a span of text from the context, similar to \citet{rajpurkar2016squad}.
We additionally ensure that it is necessary to perform multi-hop reasoning to correctly answer the question.

It is non-trivial to collect text-based multi-hop questions.
In our pilot studies, we found that simply giving an arbitrary set of paragraphs to crowd workers is counterproductive, because for most paragraph sets, it is difficult to ask a meaningful multi-hop question.
To address this challenge, we carefully design a pipeline to collect text-based multi-hop questions.
Below, we will highlight the key design choices in our pipeline.

\paragraph{Building a Wikipedia Hyperlink Graph.}
We use the entire English Wikipedia dump as our corpus.\footnote{\url{https://dumps.wikimedia.org/}}
In this corpus, we make two observations:
(1) hyper-links in the Wikipedia articles often naturally entail a relation between two (already disambiguated) entities in the context, which could potentially be used to facilitate multi-hop reasoning;
(2) the first paragraph of each article often contains much information that could be queried in a meaningful way.
Based on these observations, we extract all the hyperlinks from the first paragraphs of all Wikipedia articles.
With these hyperlinks, we build a directed graph $G$, where each edge $(a, b)$ indicates there is a hyperlink from the first paragraph of article $a$ to article $b$.

\paragraph{Generating Candidate Paragraph Pairs.}
To generate meaningful pairs of paragraphs for multi-hop question answering with $G$, we start by considering an example question ``when was the singer and songwriter of Radiohead born?''
To answer this question, one would need to first reason that the ``singer and songwriter of Radiohead'' is ``Thom Yorke'', and then figure out his birthday in the text.
We call ``Thom Yorke'' a \emph{bridge entity} in this example.
Given an edge $(a, b)$ in the hyperlink graph $G$, the entity of $b$ can usually be viewed as a bridge entity that connects $a$ and $b$.
%In other words, bridge entities $b$ create possibilities of multi-hop reasoning.
As we observe articles $b$ usually determine the theme of the shared context between $a$ and $b$, but not all articles $b$ are suitable for collecting multi-hop questions.
% However, not all entity $b$'s are suitable for collecting multi-hop question, as we observe they usually determine the theme of the shared context between $a$ and $b$.
For example, entities like countries are frequently referred to in Wikipedia, but don't necessarily have much in common with all incoming links.
It is also difficult, for instance, for the crowd workers to ask meaningful multi-hop questions about highly technical entities like the IPv4 protocol.
To alleviate this issue, we constrain the bridge entities to a set of manually curated pages in Wikipedia (see Appendix \ref{sec:collection_details}).
After curating a set of pages $B$, we create candidate paragraph pairs by sampling edges $(a, b)$ from the hyperlink graph such that $b \in B$.

\paragraph{Comparison Questions.}
In addition to questions collected using bridge entities, we also collect another type of multi-hop questions---comparison questions.
The main idea is that comparing two entities from the same category usually results in interesting multi-hop questions, e.g., ``Who has played for more NBA teams, Michael Jordan or Kobe Bryant?''
To facilitate collecting this type of question, we manually curate 42 lists of similar entities (denoted as $L$) from Wikipedia.\footnote{This is achieved by manually curating lists from the Wikipedia ``List of lists of lists'' (\url{https://wiki.sh/y8qv}). One example is ``Highest Mountains on Earth''.}
To generate candidate paragraph pairs, we randomly sample two paragraphs from the same list and present them to the crowd worker.

To increase the diversity of multi-hop questions, we also introduce a subset of yes/no questions in comparison questions.
This complements the original scope of comparison questions by offering new ways to require systems to reason over both paragraphs.
For example, consider the entities Iron Maiden (from the UK) and AC/DC (from Australia).
Questions like ``Is Iron Maiden or AC/DC from the UK?'' are not ideal, because one would deduce the answer is ``Iron Maiden'' even if one only had access to that article.
With yes/no questions, one may ask ``Are Iron Maiden and AC/DC from the same country?'', which requires reasoning over both paragraphs.

To the best of our knowledge, text-based comparison questions are a novel type of questions that have not been considered by previous datasets.
More importantly, answering these questions usually requires arithmetic comparison, such as comparing ages given birth dates, which presents a new challenge for future model development.

\paragraph{Collecting Supporting Facts.}
To enhance the explainability of question answering systems, we want them to output a set of \emph{supporting facts} necessary to arrive at the answer, when the answer is generated.
To this end, we also collect the sentences that determine the answers from crowd workers.
These supporting facts can serve as strong supervision for what sentences to pay attention to.
Moreover, we can now test the explainability of a model by comparing the predicted supporting facts to the ground truth ones.

The overall procedure of data collection is illustrated in Algorithm \ref{algo:collect}.
% We leave details about curating Wikipedia pages to Appendix \ref{sec:collection_details}.

\begin{algorithm}[t]
\small
    \caption{Overall data collection procedure}
    \label{algo:collect}
\begin{algorithmic}
    \STATE {\bfseries Input:} question type ratio $r_1=0.75$, yes/no ratio $r_2=0.5$
    \WHILE{not finished}
        \IF {$\mathrm{random()} < r_1$}
            \STATE Uniformly sample an entity $b \in B$
            \STATE Uniformly sample an edge $(a, b)$
            \STATE Workers ask a question about paragraphs $a$ and $b$
        \ELSE
            \STATE Sample a list from $L$, with probabilities weighted by list sizes
            % \todo{Aren't we doing entity uniform in the end?}
            \STATE Uniformly sample two entities $(a, b)$ from the list
            \IF {$\mathrm{random()} < r_2$}
                \STATE Workers ask a yes/no question to compare $a$ and $b$
            \ELSE
                \STATE Workers ask a question with a span answer to compare $a$ and $b$
            \ENDIF
        \ENDIF
        \STATE Workers provide the supporting facts
    \ENDWHILE
\end{algorithmic}
\end{algorithm}

% 1. Link graph, shortest path

% 2. Filtering: Type B (pick clusteroids) and Type 4 (pick lists)

% 3. UI (how are supfacts collected?)

% 4. Split

% 5. Settings and Distractors